%============================================================
%  Bd. 06 - Injektive Funktionen
%============================================================

\documentclass[main.tex]{subfiles}

\ifSubfilesClassLoaded{
  \BandSetupStandalone{B06}
}{
}

\title{Bd. 06 - Injektive Funktionen}
\author{Martin Kunze}
\date{}
\setcounter{file}{6}

\hypersetup{
  pdftitle={Bd. 06 - Injektive Funktionen},
  pdfauthor={Martin Kunze}
}

\begin{document}

\title{Bd. 06 - Injektive Funktionen}
\author{Martin Kunze}
\date{}
\hypersetup{pageanchor=false}
\maketitle
\hypersetup{pageanchor=true}
\setcounter{tocdepth}{3}
\tableofcontents
\setcounter{file}{6}
\renewcommand{\FormulaBandID}{6}

\chapter{Einleitung}

Dieser Band baut auf der allgemeinen Funktionstheorie aus Band~05 auf. Wie im
axiomatischen Aufbau von Band~03 werden zuerst Begriff und Grundaxiom
zusammengeführt; danach folgen Konsequenzen und charakteristische Beispiele.
Hier werden künftig genau die Konstruktionen gesammelt, deren maßgebliche
Eigenschaft die Injektivität ist.

\chapter{Axiomatischer Aufbau der injektiven Funktionen}

\section{Begriff und Grundaxiom}

\FormulaDefDeltaK[Begriff der injektiven Funktion]{F\colon A\rightarrowtail B}{Injektive Funktion}{
  \DeltaRow{Mengen}{A\dsep B\dsep x\dsep y}
  \DeltaPrem{Funktionen}{F\colon A\to B}[\FormulaRefAuto{Funktion}]
  %
  \DeltaRow{\textbf{Axiome}}{}
  %
  % — Injektivität —
  \DeltaRow{Injektivität}
           {F\colon A\inj B\dsep x,y\in A\dsep F(x)=F(y)\vdash x=y}
           [\FormulaRefAuto{F\colon A\inj B\dsep x\in A\dsep y\in A\dsep F(x)=F(y)\vdash x=y}]
  \DeltaRow{\textbf{Neue Symbole}}{}
  \DeltaRow{Injektive Funktionen}{}
}

\FormulaAxiomDeltaK[Injektivität]%
{F\colon A\inj B\dsep x\in A\dsep y\in A\dsep F(x)=F(y)\vdash x=y}%
{Injektivität}{
\DeltaRow{Mengen}{x\dsep y\dsep A\dsep B}
}

\FormulaThmDelta{F\colon A\inj B\dsep x\in A\vdash F(x)\in B}{
\DeltaRow{Mengen}{A\dsep B}
}
\begin{tabproof}
  \proofstep{1}{F\colon A\inj B}{\rA}
  \proofstep{2}{x\in A}{\rA}
  \proofstep{1}{F\colon A\to B}{\FormulaRefAuto{Injektive Funktion}{1}}
  \proofstep{1,2}{F(x)\in B}{%
    \FormulaRefAuto{F\colon A\to B\dsep x\in A\vdash F(x)\in B}{3,2}}
\end{tabproof}

\FormulaThmDelta[Elemente einer Teilmenge liegen im Bild der Teilmenge unter injektiven Funktionen]{%
  C\subseteq A\dsep F\colon A\inj B\dsep x\in C\vdash F(x)\in F[C]
}{
\DeltaRow{Mengen}{A\dsep B\dsep C\dsep x}
\DeltaRow{Funktionen}{F}
}
\begin{tabproof}
  \proofstep{1}{C\subseteq A}{\rA}
  \proofstep{2}{F\colon A\inj B}{\rA}
  \proofstep{3}{x\in C}{\rA}
  \proofstep{1,2,3}{F(x)\in F[C]}{%
    \FormulaRefAuto{C\subseteq A\dsep F\colon A\to B\dsep x\in C\vdash F(x)\in F[C]}{%
      1,\FormulaRefAuto{Injektive Funktion}{2},3}}
\end{tabproof}

\FormulaThmDelta[Elemente des Definitionsbereichs liegen im Bild unter injektiven Funktionen]{%
  F\colon A\inj B\dsep x\in A\vdash F(x)\in F[A]
}{
\DeltaRow{Mengen}{A\dsep B\dsep x}
\DeltaRow{Funktionen}{F}
}
\begin{tabproof}
  \proofstep{1}{F\colon A\inj B}{\rA}
  \proofstep{2}{x\in A}{\rA}
  \proofstep{}{A\subseteq A}{\FormulaRefAuto{A\subseteq A}}
  \proofstep{1,2}{F(x)\in F[A]}{%
    \FormulaRefAuto{C\subseteq A\dsep F\colon A\inj B\dsep
    x\in C\vdash F(x)\in F[C]}{3,1,2}}
\end{tabproof}


\FormulaThmDelta%
{F\colon A\inj B\dsep (x,z_1)\in F\dsep (y,z_2)\in F\dsep F(x)=F(y)\vdash x=y}%
{
\DeltaRow{Mengen}{x\dsep y\dsep z_1\dsep z_2\dsep A\dsep B}
\DeltaRow{Funktionen}{F}
}
\begin{tabproof}
  \proofstep{1}{F\colon A\inj B}{\rA}
  \proofstep{2}{(x,z_1)\in F}{\rA}
  \proofstep{3}{(y,z_2)\in F}{\rA}
  \proofstep{4}{F(x)=F(y)}{\rA}
  \proofstep{1}{F\colon A\to B}{\FormulaRefAuto{Injektive Funktion}{1}}
  \proofstep{1,2}{x\in A}{%
    \FormulaRefAuto{F\colon A\to B\dsep (x,y)\in F\vdash x\in A}{5,2}}
  \proofstep{1,3}{y\in A}{%
    \FormulaRefAuto{F\colon A\to B\dsep (x,y)\in F\vdash x\in A}{5,3}}
  \proofstep{1,2,3,4}{x=y}{%
    \FormulaRefAuto{F\colon A\inj B\dsep x\in A\dsep y\in A\dsep F(x)=F(y)\vdash x=y}{1,6,7,4}}
\end{tabproof}

\FormulaThmDelta{%
  F\colon M\inj N \dsep x\in M \dsep z\in M \dsep y=F(x) \dsep y=F(z) \vdash x=z
}{%
  \DeltaRow{Mengen}{M\dsep N\dsep x\dsep y\dsep z}
  \DeltaRow{Funktionen}{F}
}
\begin{tabproof}
  \proofstep{1}{F\colon M\inj N}{\rA}
  \proofstep{2}{x\in M}{\rA}
  \proofstep{3}{z\in M}{\rA}
  \proofstep{4}{y=F(x)}{\rA}
  \proofstep{5}{y=F(z)}{\rA}
  \proofstep{4,5}{F(x)=F(z)}{\FormulaRefAuto{a = b,\, a = c \vdash b = c}{4,5}}
  \proofstep{1,2,3,6}{x=z}{%
    \FormulaRefAuto{F\colon A\inj B\dsep x\in A\dsep y\in A\dsep F(x)=F(y)\vdash x=y}{1,2,3,6}}
\end{tabproof}


\section{Bildmengen unter injektiven Funktionen}

\FormulaThmDelta{%
  F\colon A\inj B \dsep Y\in\powerset(A)\dsep x\in A\dsep F(x)\in F[Y] \vdash x\in Y
}{%
  \DeltaRow{Mengen}{A\dsep B\dsep Y\dsep x\dsep z}
  \DeltaRow{Funktionen}{F}
}
\begin{tabproof}
  \proofstep{1}{F\colon A\inj B}{\rA}
  \proofstep{2}{Y\in\powerset(A)}{\rA}
  \proofstep{3}{x\in A}{\rA}
  \proofstep{4}{F(x)\in F[Y]}{\rA}
  \proofstep{1}{F\colon A\to B}{\FormulaRefAuto{Injektive Funktion}{1}}
  \proofstep{2}{Y\subseteq A}{\FormulaRefAuto{x \in \powerset(A)\eqvdash x \subseteq A}{2}}
  \proofstep{1,2,4}{\exists z\in Y\,F(x)=F(z)}{%
    \FormulaRefAuto{C\subseteq A\dsep F\colon A\to B\dsep y\in F[C]\vdash \exists x\in C\,y=F(x)}{6,5,4}}

  \proofstep{8}{z\in Y\land F(x)=F(z)}{\rA}
  \proofstep{8}{z\in Y}{\rAEa{8}}
  \proofstep{8}{F(x)=F(z)}{\rAEb{8}}
  \proofstep{2,8}{z\in A}{\FormulaRefAuto{A\subseteq B\dsep x\in A \vdash x\in B}{6,9}}
  \proofstep{1,2,3,8}{x=z}{\FormulaRefAuto{Injektivität}{1,3,11,10}}
  \proofstep{1,2,3,8}{z=x}{\FormulaRefAuto{a=b\vdash b=a}{12}}
  \proofstep{1,2,3,8}{x\in Y}{\rIE{13,9}}

  \proofstep{1,2,3,4}{x\in Y}{\rEE{7,8,14}}
\end{tabproof}

\FormulaThmDelta[Injektive Funktionen sind auf Potenzmengen injektiv]{%
  F\colon A\inj B \dsep X\in\powerset(A)\dsep Y\in\powerset(A)\dsep F[X]=F[Y] \vdash X=Y
}{%
  \DeltaRow{Mengen}{A\dsep B\dsep X\dsep Y\dsep x}
  \DeltaRow{Funktionen}{F}
}
\begin{tabproofsplitwide}
  \proofpartwideindR[Teilmengenrichtung \(X\subseteq Y\)]{%
    F\colon A\inj B \dsep X\in\powerset(A)\dsep Y\in\powerset(A)\dsep F[X]=F[Y] \vdash X\subseteq Y
  }{%
    F\colon A\inj B \dsep X\in\powerset(A)\dsep Y\in\powerset(A)\dsep F[X]=F[Y] \vdash X\subseteq Y
  }
    \proofstepwidestar[1]{F\colon A\inj B}{\rA}
    \proofstepwidestar[2]{X\in\powerset(A)}{\rA}
    \proofstepwidestar[3]{Y\in\powerset(A)}{\rA}
    \proofstepwidestar[4]{F[X]=F[Y]}{\rA}
    \proofstepwidestar[2]{X\subseteq A}{\FormulaRefAuto{x \in \powerset(A)\eqvdash x \subseteq A}{2}}
    \proofstepwidestar[1]{F\colon A\to B}{\FormulaRefAuto{Injektive Funktion}{1}}

    \proofstepwidestar[7]{x\in X}{\rA}
    \proofstepwidestar[2,7]{x\in A}{\FormulaRefAuto{A\subseteq B\dsep x\in A \vdash x\in B}{5,7}}
    \proofstepwidestar[1,2,7]{F(x)\in F[X]}{%
      \FormulaRefAuto{C\subseteq A\dsep F\colon A\to B\dsep x\in C\vdash F(x)\in F[C]}{5,6,7}}
    \proofstepwidestar[1,2,3,4,7]{F(x)\in F[Y]}{\rIE{4,9}}
    \proofstepwidestar[1,2,3,4,7]{x\in Y}{%
      \FormulaRefAuto{F\colon A\inj B \dsep Y\in\powerset(A)\dsep x\in A\dsep F(x)\in F[Y] \vdash x\in Y}{1,3,8,10}}

    \proofstepwidestar[1,2,3,4]{x\in X\rightarrow x\in Y}{\rRI{7,11}}
    \proofstepwidestar[1,2,3,4]{\forall x\,(x\in X\rightarrow x\in Y)}{\rUI{\rRI{7,12}}}
    \proofstepwidestar[1,2,3,4]{X\subseteq Y}{%
      \FormulaRefAuto{A\subseteq B \coloneq \forall x\,(x\in A \rightarrow x\in B)}{13}}
  \closeproofpartwideindR

  \proofpartwideindR[Teilmengenrichtung \(Y\subseteq X\)]{%
    F\colon A\inj B \dsep X\in\powerset(A)\dsep Y\in\powerset(A)\dsep F[X]=F[Y] \vdash Y\subseteq X
  }{%
    F\colon A\inj B \dsep X\in\powerset(A)\dsep Y\in\powerset(A)\dsep F[X]=F[Y] \vdash Y\subseteq X
  }
    \proofstepwidestar[1]{F\colon A\inj B}{\rA}
    \proofstepwidestar[2]{X\in\powerset(A)}{\rA}
    \proofstepwidestar[3]{Y\in\powerset(A)}{\rA}
    \proofstepwidestar[4]{F[X]=F[Y]}{\rA}
    \proofstepwidestar[4]{F[Y]=F[X]}{\FormulaRefAuto{a=b\vdash b=a}{4}}
    \proofstepwidestar[1,2,3,4]{Y\subseteq X}{%
      \FormulaRefAuto{F\colon A\inj B \dsep X\in\powerset(A)\dsep Y\in\powerset(A)\dsep F[X]=F[Y] \vdash X\subseteq Y}{1,3,2,5}}
  \closeproofpartwideindR

  \proofpartwide{Schluss}
    \proofstepwidestar[1]{F\colon A\inj B}{\rA}
    \proofstepwidestar[2]{X\in\powerset(A)}{\rA}
    \proofstepwidestar[3]{Y\in\powerset(A)}{\rA}
    \proofstepwidestar[4]{F[X]=F[Y]}{\rA}
    \proofstepwidestar[1,2,3,4]{X\subseteq Y}{%
      \FormulaRefAuto{F\colon A\inj B \dsep X\in\powerset(A)\dsep Y\in\powerset(A)\dsep F[X]=F[Y] \vdash X\subseteq Y}{1,2,3,4}}
    \proofstepwidestar[1,2,3,4]{Y\subseteq X}{%
      \FormulaRefAuto{F\colon A\inj B \dsep X\in\powerset(A)\dsep Y\in\powerset(A)\dsep F[X]=F[Y] \vdash Y\subseteq X}{1,2,3,4}}
    \proofstepwidestar[1,2,3,4]{X=Y}{%
      \FormulaRefAuto{A \subseteq B, B \subseteq A \vdash A = B}{5,6}}
  \closeproofpartwide
\end{tabproofsplitwide}

\FormulaThmDelta[Zeugen gemeinsamer Bildwerte liegen im Durchschnitt]{%
  F\colon M\inj N \dsep A\subseteq M \dsep B\subseteq M \dsep x\in A \dsep z\in B \dsep y=F(x) \dsep y=F(z) \vdash x\in A\cap B
}{%
  \DeltaRow{Mengen}{A\dsep B\dsep M\dsep N\dsep x\dsep y\dsep z}
  \DeltaRow{Funktionen}{F}
}
\begin{tabproof}
  \proofstep{1}{F\colon M\inj N}{\rA}
  \proofstep{2}{A\subseteq M}{\rA}
  \proofstep{3}{B\subseteq M}{\rA}
  \proofstep{4}{x\in A}{\rA}
  \proofstep{5}{z\in B}{\rA}
  \proofstep{6}{y=F(x)}{\rA}
  \proofstep{7}{y=F(z)}{\rA}
  \proofstep{2,4}{x\in M}{\FormulaRefAuto{A\subseteq B\dsep x\in A\vdash x\in B}{2,4}}
  \proofstep{3,5}{z\in M}{\FormulaRefAuto{A\subseteq B\dsep x\in A\vdash x\in B}{3,5}}
  \proofstep{1,6,7,8,9}{x=z}{%
    \FormulaRefAuto{F\colon M\inj N \dsep x\in M \dsep z\in M \dsep y=F(x) \dsep y=F(z) \vdash x=z}{1,8,9,6,7}}
  \proofstep{5,10}{x\in B}{\rIE{10,5}}
  \proofstep{4,11}{x\in A\land x\in B}{\rAI{4,11}}
  \proofstep{1,2,3,4,5,6,7}{x\in A\cap B}{%
    \FormulaRefAuto{x\in A\cap B\eqvdash x\in A\land x\in B}{12}}
\end{tabproof}


\FormulaThmDelta[Gemeinsame Bildwerte liegen im Bild des Durchschnitts]{%
  F\colon M\inj N \dsep A\subseteq M \dsep B\subseteq M \dsep y\in F[A] \dsep y\in F[B] \vdash y\in F[A\cap B]
}{%
  \DeltaRow{Mengen}{A\dsep B\dsep M\dsep N\dsep x\dsep y\dsep z}
  \DeltaRow{Funktionen}{F}
}
\begin{tabproof}
  \proofstep{1}{F\colon M\inj N}{\rA}
  \proofstep{2}{A\subseteq M}{\rA}
  \proofstep{3}{B\subseteq M}{\rA}
  \proofstep{4}{y\in F[A]}{\rA}
  \proofstep{5}{y\in F[B]}{\rA}
  \proofstep{1}{F\colon M\to N}{\FormulaRefAuto{Injektive Funktion}{1}}
  \proofstep{2,6,4}{\exists x\in A\,y=F(x)}{%
    \FormulaRefAuto{C\subseteq A\dsep F\colon A\to B\dsep y\in F[C]\vdash \exists x\in C\,y=F(x)}{2,6,4}}
  \proofstep{3,6,5}{\exists z\in B\,y=F(z)}{%
    \FormulaRefAuto{C\subseteq A\dsep F\colon A\to B\dsep y\in F[C]\vdash \exists x\in C\,y=F(x)}{3,6,5}}
  \proofstep{9}{x\in A\land y=F(x)}{\rA}
  \proofstep{10}{z\in B\land y=F(z)}{\rA}
  \proofstep{1,2,3,9,10}{x\in A\cap B}{%
    \FormulaRefAuto{F\colon M\inj N \dsep A\subseteq M \dsep B\subseteq M \dsep x\in A \dsep z\in B \dsep y=F(x) \dsep y=F(z) \vdash x\in A\cap B}{%
      1,2,3,\rAEa{9},\rAEa{10},\rAEb{9},\rAEb{10}}}
  \proofstep{1,2,3,9,10}{\exists x\in A\cap B\,y=F(x)}{%
    \rEI{\rAI{11,\rAEb{9}}}}
  \proofstep{1,2,3,9,10}{y\in F[A\cap B]}{%
    \FormulaRefAuto{C\subseteq A\dsep F\colon A\to B\dsep \exists x\in C\,y=F(x)\vdash y\in F[C]}{%
      \FormulaRefAuto{A\subseteq B,\, B\subseteq C \vdash A\subseteq C}{%
        \FormulaRefAuto{A\cap B\subseteq A},2},%
      6,12}}
  \proofstep{1,2,3,5,9}{y\in F[A\cap B]}{\rEE{8,10,13}}
  \proofstep{1,2,3,4,5}{y\in F[A\cap B]}{\rEE{7,9,14}}
\end{tabproof}


\FormulaThmDelta[Bild eines Durchschnitts bei injektiven Funktionen]{%
  F\colon M\inj N \dsep A\subseteq M \dsep B\subseteq M \vdash F[A\cap B]=F[A]\cap F[B]
}{%
  \DeltaRow{Mengen}{A\dsep B\dsep M\dsep N\dsep x\dsep y\dsep z}
  \DeltaRow{Funktionen}{F}
}
\begin{tabproofsplitwide}

  \proofpartwideind[Teilmengenrichtung \(F[A\cap B]\subseteq F[A]\cap F[B]\)]{%
    F\colon M\inj N \dsep A\subseteq M \dsep B\subseteq M \vdash F[A\cap B]\subseteq F[A]\cap F[B]
  }[%
    F\colon M\inj N \dsep A\subseteq M \dsep B\subseteq M \vdash F[A\cap B]\subseteq F[A]\cap F[B]
  ]
    \proofstepwidestar[1]{F\colon M\inj N}{\rA}
    \proofstepwidestar[2]{A\subseteq M}{\rA}
    \proofstepwidestar[3]{B\subseteq M}{\rA}
    \proofstepwidestar[1]{F\colon M\to N}{\FormulaRefAuto{Injektive Funktion}{1}}
    \proofstepwidestar[1,2,3]{F[A\cap B]\subseteq F[A]\cap F[B]}{%
      \FormulaRefAuto{F\colon M\to N \dsep A\subseteq M \dsep B\subseteq M \vdash F[A\cap B]\subseteq F[A]\cap F[B]}{4,2,3}}
  \closeproofpartwideind

   \proofpartwideind[Teilmengenrichtung \(F[A]\cap F[B]\subseteq F[A\cap B]\)]{%
    F\colon M\inj N \dsep A\subseteq M \dsep B\subseteq M \vdash F[A]\cap F[B]\subseteq F[A\cap B]
  }[%
    F\colon M\inj N \dsep A\subseteq M \dsep B\subseteq M \vdash F[A]\cap F[B]\subseteq F[A\cap B]
  ]
    \proofstepwidestar[1]{F\colon M\inj N}{\rA}
    \proofstepwidestar[2]{A\subseteq M}{\rA}
    \proofstepwidestar[3]{B\subseteq M}{\rA}
    \proofstepwidestar[4]{y\in F[A]\cap F[B]}{\rA}
    \proofstepwidestar[4]{y\in F[A]\land y\in F[B]}{%
      \FormulaRefAuto{y\in A\cap B\eqvdash y\in A\land y\in B}{4}}
    \proofstepwidestar[4]{y\in F[A]}{\rAEa{5}}
    \proofstepwidestar[4]{y\in F[B]}{\rAEb{5}}
    \proofstepwidestar[1,2,3,6,7]{y\in F[A\cap B]}{%
      \FormulaRefAuto{F\colon M\inj N \dsep A\subseteq M \dsep B\subseteq M \dsep y\in F[A] \dsep y\in F[B] \vdash y\in F[A\cap B]}{1,2,3,6,7}}
    \proofstepwidestar[1]{\forall y\,\bigl(y\in F[A]\cap F[B]\rightarrow y\in F[A\cap B]\bigr)}{\rUI{\rRI{4,8}}}
    \proofstepwidestar[1]{F[A]\cap F[B]\subseteq F[A\cap B]}{%
      \FormulaRefAuto{A\subseteq B \coloneq \forall x\,(x\in A \rightarrow x\in B)}{9}}
  \closeproofpartwideind

  \proofpartwide{Zusammenfassung}
    \proofstepwidestar[1]{F\colon M\inj N}{\rA}
    \proofstepwidestar[2]{A\subseteq M}{\rA}
    \proofstepwidestar[3]{B\subseteq M}{\rA}
    \proofstepwidestar[1,2,3]{F[A\cap B]\subseteq F[A]\cap F[B]}{%
      \FormulaRefAuto{F\colon M\inj N \dsep A\subseteq M \dsep B\subseteq M \vdash F[A\cap B]\subseteq F[A]\cap F[B]}{1,2,3}}
    \proofstepwidestar[1,2,3]{F[A]\cap F[B]\subseteq F[A\cap B]}{%
      \FormulaRefAuto{F\colon M\inj N \dsep A\subseteq M \dsep B\subseteq M \vdash F[A]\cap F[B]\subseteq F[A\cap B]}{1,2,3}}
    \proofstepwidestar[1,2,3]{F[A\cap B]=F[A]\cap F[B]}{%
      \FormulaRefAuto{A\subseteq B\dsep B\subseteq A \vdash A=B}{4,5}}
  \closeproofpartwide

\end{tabproofsplitwide}



\section{Urbildmengen unter injektiven Funktionen}

\FormulaThmDelta[Urbilder liegen im Definitionsbereich injektiver Funktionen]{%
  C\subseteq B\dsep F\colon A\inj B \vdash F^{-1}[C]\subseteq A
}{
  \DeltaRow{Mengen}{A\dsep B\dsep C}
  \DeltaRow{Funktionen}{F}
}
\begin{tabproof}
  \proofstep{1}{C\subseteq B}{\rA}
  \proofstep{2}{F\colon A\inj B}{\rA}
  \proofstep{2}{F\colon A\to B}{\FormulaRefAuto{Injektive Funktion}{2}}
  \proofstep{1,2}{F^{-1}[C]\subseteq A}{%
    \FormulaRefAuto{F\colon A\to B\dsep C\subseteq B \vdash F^{-1}[C]\subseteq A}{3,1}}
\end{tabproof}


\chapter{Beispiele und Konstruktionen}

\section{Zielmengenwechsel und Einschränkungen}

\subsection{Zielmengenwechsel}

\FormulaThmDelta[Verkleinerung der Zielmenge erhält Injektivität]{%
  F\colon A\inj B\dsep \forall x\in A\,F(x)\in C
  \vdash
  F\colon A\inj C
}{%
  \DeltaDecl{Mengen}{A\dsep B\dsep C}%
  \DeltaDecl{Funktionen}{F}%
}
\begin{tabproof}
  \proofstep{1}{F\colon A\inj B}{\rA}
  \proofstep{2}{\forall x\in A\,F(x)\in C}{\rA}
  \proofstep{1}{F\colon A\to B}{\FormulaRefAuto{Injektive Funktion}{1}}
  \proofstep{1,2}{F\colon A\to C}{\FormulaRefAuto{F\colon A\to B\dsep \forall x\in A\,F(x)\in C\vdash F\colon A\to C}{3,2}}
  \proofstep{1}{x\in A\land y\in A\land F(x)=F(y)\rightarrow x=y}{%
    \FormulaRefAuto{F\colon A\inj B\dsep x\in A\dsep y\in A\dsep F(x)=F(y)\vdash x=y}{1}}
  \proofstep{1,2}{F\colon A\inj C}{\FormulaRefAuto{Injektive Funktion}{4,5}}
\end{tabproof}


\FormulaThmDeltaK[Erweiterung der Zielmenge erhält Injektivität]{%
  F\colon A\inj B\dsep B\subseteq C\vdash F\colon A\inj C
}{InjectiveCodomainExtension}{%
  \DeltaDecl{Mengen}{A\dsep B\dsep C}%
  \DeltaDecl{Funktionen}{F}%
}
\begin{tabproof}
  \proofstep{1}{F\colon A\inj B}{\rA}
  \proofstep{2}{B\subseteq C}{\rA}
  \proofstep{1}{F\colon A\to B}{\FormulaRefAuto{Injektive Funktion}{1}}
  \proofstep{1,2}{F\colon A\to C}{%
    \FormulaRefAuto{F\colon A\to B\dsep B\subseteq C\vdash F\colon A\to C}{3,2}}
  \proofstep{1}{x\in A\land y\in A\land F(x)=F(y)\rightarrow x=y}{%
    \FormulaRefAuto{F\colon A\inj B\dsep x\in A\dsep y\in A\dsep
    F(x)=F(y)\vdash x=y}{1}}
  \proofstep{1,2}{F\colon A\inj C}{\FormulaRefAuto{Injektive Funktion}{4,5}}
\end{tabproof}

\subsection{Einschränkungen}

\subsection{Erhalt von Eigenschaften}

\FormulaThmDelta[Injektivität der Einschränkung]{%
F\colon A\inj B\dsep C\subseteq A \dsep x\in C \dsep y\in C \dsep
F\restriction_C(x)=F\restriction_C(y) \vdash x=y
}{
  \DeltaRow{Mengen}{A\dsep B\dsep C\dsep x\dsep y}
  \DeltaRow{Funktionen}{F}
  \DeltaPrem{Funktionen}{F\restriction_C\colon C\to B}[\FormulaRefAuto{C\subseteq A \vdash F\restriction_C\colon C\to B}]
}
\begin{tabproof}
  \proofstep{1}{F\colon A\inj B}{\rA}
  \proofstep{2}{C\subseteq A}{\rA}
  \proofstep{3}{x\in C}{\rA}
  \proofstep{4}{y\in C}{\rA}
  \proofstep{5}{F\restriction_C(x)=F\restriction_C(y)}{\rA}
  \proofstep{2,3,4,5}{F(x)=F(y)}{%
    \FormulaRefAuto{C\subseteq A \dsep x\in C \dsep y\in C \dsep F\restriction_C(x)=F\restriction_C(y)\vdash F(x)=F(y)}{2,3,4,5}}
  \proofstep{2,3}{x\in A}{\FormulaRefAuto{A\subseteq B, x\in A\vdash x\in B}{2,3}}
  \proofstep{2,4}{y\in A}{\FormulaRefAuto{A\subseteq B, x\in A\vdash x\in B}{2,4}}
  \proofstep{1,2,3,4,5}{x=y}{%
    \FormulaRefAuto{F\colon A\inj B\dsep x\in A\dsep y\in A\dsep F(x)=F(y)\vdash x=y}{1,7,8,6}}
\end{tabproof}

\FormulaThmDelta[Einschränkung einer injektiven Funktion]{%
 F\colon A\inj B\dsep C\subseteq A \vdash F\restriction_C\colon C\inj B
}{
  \DeltaRow{Mengen}{A\dsep B\dsep C}
  \DeltaRow{Funktionen}{F}
}
\begin{tabproof}
  \proofstep{1}{F\colon A\inj B}{\rA}
  \proofstep{2}{C\subseteq A}{\rA}
  \proofstep{2}{F\restriction_C\colon C\to B}{%
    \FormulaRefAuto{C\subseteq A \vdash F\restriction_C\colon C\to B}{2}}
  \proofstep{1,2}{\forall x,y\in C\,(F\restriction_C(x)=F\restriction_C(y)\rightarrow x=y)}{%
    \FormulaRefAuto{F\colon A\inj B\dsep C\subseteq A \dsep x\in C \dsep y\in C \dsep
F\restriction_C(x)=F\restriction_C(y) \vdash x=y}{1,2}}
  \proofstep{1,2}{F\restriction_C\colon C\inj B}{%
    \FormulaRefAuto{Injektive Funktion}{3,4}}
\end{tabproof}


\subsection{Einschränkung auf Urbilder}

\FormulaThmDelta[Einschränkung auf ein Urbild als injektive Funktion]{%
  F\colon A\inj B\dsep T\subseteq B \vdash F\restriction_{F^{-1}[T]}\colon F^{-1}[T]\inj T
}{
  \DeltaRow{Mengen}{A\dsep B\dsep T\dsep x}
  \DeltaRow{Funktionen}{F}
}
\begin{tabproof}
  \proofstep{1}{F\colon A\inj B}{\rA}
  \proofstep{2}{T\subseteq B}{\rA}
  \proofstep{1}{F\colon A\to B}{\FormulaRefAuto{Injektive Funktion}{1}}
  \proofstep{2,3}{F^{-1}[T]\subseteq A}{%
    \FormulaRefAuto{F\colon A\to B\dsep C\subseteq B \vdash F^{-1}[C]\subseteq A}{3,2}}
  \proofstep{1,4}{F\restriction_{F^{-1}[T]}\colon F^{-1}[T]\inj B}{%
    \FormulaRefAuto{F\colon A\inj B\dsep C\subseteq A \vdash F\restriction_C\colon C\inj B}{1,4}}
  \proofstep{2,3}{F\restriction_{F^{-1}[T]}\colon F^{-1}[T]\to T}{%
    \FormulaRefAuto{F\colon A\to B\dsep T\subseteq B \vdash F\restriction_{F^{-1}[T]}\colon F^{-1}[T]\to T}{3,2}}
  \proofstep{7}{x\in F^{-1}[T]}{\rA}
  \proofstep{2,6,7}{F\restriction_{F^{-1}[T]}(x)\in T}{%
    \FormulaRefAuto{F\colon A\to B\dsep x\in A\vdash F(x)\in B}{6,7}}
  \proofstep{2,6}{\forall x\in F^{-1}[T]\,F\restriction_{F^{-1}[T]}(x)\in T}{%
    \rUI{\rRI{7,8}}}
  \proofstep{1,2}{F\restriction_{F^{-1}[T]}\colon F^{-1}[T]\inj T}{%
    \FormulaRefAuto{F\colon A\inj B\dsep \forall x\in A\,F(x)\in C \vdash F\colon A\inj C}{5,9}}
\end{tabproof}

\section{Fallunterscheidungsabbildungen}

\FormulaThmDeltaK[Injektives Zusammenführen von Fallabbildungszweigen]{%
  A\cap B=\varnothing\dsep
  f\colon A\inj C\dsep g\colon B\inj C\dsep
  f[A]\cap g[B]=\varnothing
  \vdash \CaseFun{A}{B}{f}{g}\colon A\cup B\inj C
}{CaseFunInjectiveDisjointImages}{%
  \DeltaRow{Mengen}{A\dsep B\dsep C\dsep x\dsep y}
  \DeltaRow{Funktionen}{f\dsep g\dsep \CaseFun{A}{B}{f}{g}}
}
\begin{tabproofsplitwide}
  \proofpartwideindR[Fall \(x\in A,\,y\in A\)]{%
    \begin{aligned}[t]
    &A\cap B=\varnothing\dsep
    f\colon A\inj C\dsep g\colon B\inj C\\
    &f[A]\cap g[B]=\varnothing\dsep
    \CaseFun{A}{B}{f}{g}(x)=\CaseFun{A}{B}{f}{g}(y)\\
    &x\in A\dsep y\in A\vdash x=y
    \end{aligned}
  }{%
    A\cap B=\varnothing\dsep
    f\colon A\inj C\dsep g\colon B\inj C\dsep
    f[A]\cap g[B]=\varnothing\dsep
    \CaseFun{A}{B}{f}{g}(x)=\CaseFun{A}{B}{f}{g}(y)\dsep
    x\in A\dsep y\in A
    \vdash x=y
  }[CaseFunInjectiveDisjointImagesCaseAA]
    \proofstepwidestar[1]{A\cap B=\varnothing}{\rA}
    \proofstepwidestar[2]{f\colon A\inj C}{\rA}
    \proofstepwidestar[3]{g\colon B\inj C}{\rA}
    \proofstepwidestar[4]{f[A]\cap g[B]=\varnothing}{\rA}
    \proofstepwidestar[5]{%
      \CaseFun{A}{B}{f}{g}(x)=\CaseFun{A}{B}{f}{g}(y)}{\rA}
    \proofstepwidestar[6]{x\in A}{\rA}
    \proofstepwidestar[7]{y\in A}{\rA}
    \proofstepwidestar[2]{f\colon A\to C}{%
      \FormulaRefAuto{Injektive Funktion}{2}}
    \proofstepwidestar[3]{g\colon B\to C}{%
      \FormulaRefAuto{Injektive Funktion}{3}}
    \proofstepwidestar[2,3,6]{\CaseFun{A}{B}{f}{g}(x)=f(x)}{%
      \FormulaRefAuto{f\colon A\to C\dsep g\colon B\to C\dsep x\in A
      \vdash \CaseFun{A}{B}{f}{g}(x)=f(x)}{8,9,6}}
    \proofstepwidestar[2,3,7]{\CaseFun{A}{B}{f}{g}(y)=f(y)}{%
      \FormulaRefAuto{f\colon A\to C\dsep g\colon B\to C\dsep x\in A
      \vdash \CaseFun{A}{B}{f}{g}(x)=f(x)}{8,9,7}}
    \proofstepwidestar[2,3,5,6,7]{f(x)=f(y)}{%
      \FormulaRefAuto{a = b\dsep a = c\dsep b=d\vdash c = d}{5,10,11}}
    \proofstepwidestar[1,2,3,4,5,6,7]{x=y}{%
      \FormulaRefAuto{F\colon A\inj B\dsep x\in A\dsep y\in A\dsep
      F(x)=F(y)\vdash x=y}{2,6,7,12}}
  \closeproofpartwideindR

  \proofpartwideindR[Fall \(x\in A,\,y\in B\)]{%
    \begin{aligned}[t]
    &A\cap B=\varnothing\dsep
    f\colon A\inj C\dsep g\colon B\inj C\\
    &f[A]\cap g[B]=\varnothing\dsep
    \CaseFun{A}{B}{f}{g}(x)=\CaseFun{A}{B}{f}{g}(y)\\
    &x\in A\dsep y\in B\vdash x=y
    \end{aligned}
  }{%
    A\cap B=\varnothing\dsep
    f\colon A\inj C\dsep g\colon B\inj C\dsep
    f[A]\cap g[B]=\varnothing\dsep
    \CaseFun{A}{B}{f}{g}(x)=\CaseFun{A}{B}{f}{g}(y)\dsep
    x\in A\dsep y\in B
    \vdash x=y
  }[CaseFunInjectiveDisjointImagesCaseAB]
    \proofstepwidestar[1]{A\cap B=\varnothing}{\rA}
    \proofstepwidestar[2]{f\colon A\inj C}{\rA}
    \proofstepwidestar[3]{g\colon B\inj C}{\rA}
    \proofstepwidestar[4]{f[A]\cap g[B]=\varnothing}{\rA}
    \proofstepwidestar[5]{%
      \CaseFun{A}{B}{f}{g}(x)=\CaseFun{A}{B}{f}{g}(y)}{\rA}
    \proofstepwidestar[6]{x\in A}{\rA}
    \proofstepwidestar[7]{y\in B}{\rA}
    \proofstepwidestar[2]{f\colon A\to C}{%
      \FormulaRefAuto{Injektive Funktion}{2}}
    \proofstepwidestar[3]{g\colon B\to C}{%
      \FormulaRefAuto{Injektive Funktion}{3}}
    \proofstepwidestar[1,2,3,7]{\CaseFun{A}{B}{f}{g}(y)=g(y)}{%
      \FormulaRefAuto{A\cap B=\varnothing\dsep f\colon A\to C\dsep
      g\colon B\to C\dsep x\in B
      \vdash \CaseFun{A}{B}{f}{g}(x)=g(x)}{1,8,9,7}}
    \proofstepwidestar[2,3,6]{\CaseFun{A}{B}{f}{g}(x)=f(x)}{%
      \FormulaRefAuto{f\colon A\to C\dsep g\colon B\to C\dsep x\in A
      \vdash \CaseFun{A}{B}{f}{g}(x)=f(x)}{8,9,6}}
    \proofstepwidestar[1,2,3,5,6,7]{f(x)=g(y)}{%
      \FormulaRefAuto{a = b\dsep a = c\dsep b=d\vdash c = d}{5,11,10}}
    \proofstepwidestar[2,6]{f(x)\in f[A]}{%
      \FormulaRefAuto{F\colon A\inj B\dsep x\in A\vdash F(x)\in F[A]}{2,6}}
    \proofstepwidestar[3,7]{g(y)\in g[B]}{%
      \FormulaRefAuto{F\colon A\inj B\dsep x\in A\vdash F(x)\in F[A]}{3,7}}
    \proofstepwidestar[1,2,3,5,6,7]{g(y)\in f[A]}{\rIE{12,13}}
    \proofstepwidestar[1,2,3,4,5,6,7]{g(y)\notin g[B]}{%
      \FormulaRefAuto{A \cap B = \varnothing,\ x \in A \vdash x \notin B}{4,15}}
    \proofstepwidestar[1,2,3,4,5,6,7]{x=y}{%
      \FormulaRefAuto{P\dsep \neg P \vdash Q}{14,16}}
  \closeproofpartwideindR

  \proofpartwideindR[Fall \(x\in B,\,y\in A\)]{%
    \begin{aligned}[t]
    &A\cap B=\varnothing\dsep
    f\colon A\inj C\dsep g\colon B\inj C\\
    &f[A]\cap g[B]=\varnothing\dsep
    \CaseFun{A}{B}{f}{g}(x)=\CaseFun{A}{B}{f}{g}(y)\\
    &x\in B\dsep y\in A\vdash x=y
    \end{aligned}
  }{%
    A\cap B=\varnothing\dsep
    f\colon A\inj C\dsep g\colon B\inj C\dsep
    f[A]\cap g[B]=\varnothing\dsep
    \CaseFun{A}{B}{f}{g}(x)=\CaseFun{A}{B}{f}{g}(y)\dsep
    x\in B\dsep y\in A
    \vdash x=y
  }[CaseFunInjectiveDisjointImagesCaseBA]
    \proofstepwidestar[1]{A\cap B=\varnothing}{\rA}
    \proofstepwidestar[2]{f\colon A\inj C}{\rA}
    \proofstepwidestar[3]{g\colon B\inj C}{\rA}
    \proofstepwidestar[4]{f[A]\cap g[B]=\varnothing}{\rA}
    \proofstepwidestar[5]{%
      \CaseFun{A}{B}{f}{g}(x)=\CaseFun{A}{B}{f}{g}(y)}{\rA}
    \proofstepwidestar[6]{x\in B}{\rA}
    \proofstepwidestar[7]{y\in A}{\rA}
    \proofstepwidestar[5]{%
      \CaseFun{A}{B}{f}{g}(y)=\CaseFun{A}{B}{f}{g}(x)}{%
      \FormulaRefAuto{a=b\vdash b=a}{5}}
    \proofstepwidestar[1,2,3,4,5,6,7]{y=x}{%
      \FormulaRefAuto{CaseFunInjectiveDisjointImagesCaseAB}{1,2,3,4,8,7,6}}
    \proofstepwidestar[1,2,3,4,5,6,7]{x=y}{%
      \FormulaRefAuto{a=b\vdash b=a}{9}}
  \closeproofpartwideindR

  \proofpartwideindR[Fall \(x\in B,\,y\in B\)]{%
    \begin{aligned}[t]
    &A\cap B=\varnothing\dsep
    f\colon A\inj C\dsep g\colon B\inj C\\
    &f[A]\cap g[B]=\varnothing\dsep
    \CaseFun{A}{B}{f}{g}(x)=\CaseFun{A}{B}{f}{g}(y)\\
    &x\in B\dsep y\in B\vdash x=y
    \end{aligned}
  }{%
    A\cap B=\varnothing\dsep
    f\colon A\inj C\dsep g\colon B\inj C\dsep
    f[A]\cap g[B]=\varnothing\dsep
    \CaseFun{A}{B}{f}{g}(x)=\CaseFun{A}{B}{f}{g}(y)\dsep
    x\in B\dsep y\in B
    \vdash x=y
  }[CaseFunInjectiveDisjointImagesCaseBB]
    \proofstepwidestar[1]{A\cap B=\varnothing}{\rA}
    \proofstepwidestar[2]{f\colon A\inj C}{\rA}
    \proofstepwidestar[3]{g\colon B\inj C}{\rA}
    \proofstepwidestar[4]{f[A]\cap g[B]=\varnothing}{\rA}
    \proofstepwidestar[5]{%
      \CaseFun{A}{B}{f}{g}(x)=\CaseFun{A}{B}{f}{g}(y)}{\rA}
    \proofstepwidestar[6]{x\in B}{\rA}
    \proofstepwidestar[7]{y\in B}{\rA}
    \proofstepwidestar[2]{f\colon A\to C}{%
      \FormulaRefAuto{Injektive Funktion}{2}}
    \proofstepwidestar[3]{g\colon B\to C}{%
      \FormulaRefAuto{Injektive Funktion}{3}}
    \proofstepwidestar[1,2,3,6]{\CaseFun{A}{B}{f}{g}(x)=g(x)}{%
      \FormulaRefAuto{A\cap B=\varnothing\dsep f\colon A\to C\dsep
      g\colon B\to C\dsep x\in B
      \vdash \CaseFun{A}{B}{f}{g}(x)=g(x)}{1,8,9,6}}
    \proofstepwidestar[1,2,3,7]{\CaseFun{A}{B}{f}{g}(y)=g(y)}{%
      \FormulaRefAuto{A\cap B=\varnothing\dsep f\colon A\to C\dsep
      g\colon B\to C\dsep x\in B
      \vdash \CaseFun{A}{B}{f}{g}(x)=g(x)}{1,8,9,7}}
    \proofstepwidestar[1,2,3,5,6,7]{g(x)=g(y)}{%
      \FormulaRefAuto{a = b\dsep a = c\dsep b=d\vdash c = d}{5,10,11}}
    \proofstepwidestar[1,2,3,4,5,6,7]{x=y}{%
      \FormulaRefAuto{F\colon A\inj B\dsep x\in A\dsep y\in A\dsep
      F(x)=F(y)\vdash x=y}{3,6,7,12}}
  \closeproofpartwideindR

  \proofpartwideindR[Fall \(x\in A\)]{%
    \begin{aligned}[t]
      &A\cap B=\varnothing\dsep f\colon A\inj C\dsep g\colon B\inj C\dsep
      f[A]\cap g[B]=\varnothing\dsep{}\\
      &\CaseFun{A}{B}{f}{g}(x)=\CaseFun{A}{B}{f}{g}(y)\dsep
      x\in A\dsep y\in A\cup B\vdash x=y
    \end{aligned}%
  }{%
    A\cap B=\varnothing\dsep f\colon A\inj C\dsep g\colon B\inj C\dsep
    f[A]\cap g[B]=\varnothing\dsep
    \CaseFun{A}{B}{f}{g}(x)=\CaseFun{A}{B}{f}{g}(y)\dsep
    x\in A\dsep y\in A\cup B\vdash x=y
  }[CaseFunInjectiveDisjointImagesCaseA]
    \proofstepwidestar[1]{A\cap B=\varnothing}{\rA}
    \proofstepwidestar[2]{f\colon A\inj C}{\rA}
    \proofstepwidestar[3]{g\colon B\inj C}{\rA}
    \proofstepwidestar[4]{f[A]\cap g[B]=\varnothing}{\rA}
    \proofstepwidestar[5]{%
      \CaseFun{A}{B}{f}{g}(x)=\CaseFun{A}{B}{f}{g}(y)}{\rA}
    \proofstepwidestar[6]{x\in A}{\rA}
    \proofstepwidestar[7]{y\in A\cup B}{\rA}
    \proofstepwidestar[7]{y\in A\lor y\in B}{%
      \FormulaRefAuto{z \in A \cup B \eqvdash z \in A \lor z \in B}{7}}
    \proofstepwidestar[9]{y\in A}{\rA}
    \proofstepwidestar[1,2,3,4,5,6,9]{x=y}{%
      \FormulaRefAuto{CaseFunInjectiveDisjointImagesCaseAA}{1,2,3,4,5,6,9}}
    \proofstepwidestar[11]{y\in B}{\rA}
    \proofstepwidestar[1,2,3,4,5,6,11]{x=y}{%
      \FormulaRefAuto{CaseFunInjectiveDisjointImagesCaseAB}{1,2,3,4,5,6,11}}
    \proofstepwidestar[1,2,3,4,5,6,7]{x=y}{\rOE{8,9,10,11,12}}
  \closeproofpartwideindR

  \proofpartwideindR[Fall \(x\in B\)]{%
    \begin{aligned}[t]
      &A\cap B=\varnothing\dsep f\colon A\inj C\dsep g\colon B\inj C\dsep
      f[A]\cap g[B]=\varnothing\dsep{}\\
      &\CaseFun{A}{B}{f}{g}(x)=\CaseFun{A}{B}{f}{g}(y)\dsep
      x\in B\dsep y\in A\cup B\vdash x=y
    \end{aligned}%
  }{%
    A\cap B=\varnothing\dsep f\colon A\inj C\dsep g\colon B\inj C\dsep
    f[A]\cap g[B]=\varnothing\dsep
    \CaseFun{A}{B}{f}{g}(x)=\CaseFun{A}{B}{f}{g}(y)\dsep
    x\in B\dsep y\in A\cup B\vdash x=y
  }[CaseFunInjectiveDisjointImagesCaseB]
    \proofstepwidestar[1]{A\cap B=\varnothing}{\rA}
    \proofstepwidestar[2]{f\colon A\inj C}{\rA}
    \proofstepwidestar[3]{g\colon B\inj C}{\rA}
    \proofstepwidestar[4]{f[A]\cap g[B]=\varnothing}{\rA}
    \proofstepwidestar[5]{%
      \CaseFun{A}{B}{f}{g}(x)=\CaseFun{A}{B}{f}{g}(y)}{\rA}
    \proofstepwidestar[6]{x\in B}{\rA}
    \proofstepwidestar[7]{y\in A\cup B}{\rA}
    \proofstepwidestar[7]{y\in A\lor y\in B}{%
      \FormulaRefAuto{z \in A \cup B \eqvdash z \in A \lor z \in B}{7}}
    \proofstepwidestar[9]{y\in A}{\rA}
    \proofstepwidestar[1,2,3,4,5,6,9]{x=y}{%
      \FormulaRefAuto{CaseFunInjectiveDisjointImagesCaseBA}{1,2,3,4,5,6,9}}
    \proofstepwidestar[11]{y\in B}{\rA}
    \proofstepwidestar[1,2,3,4,5,6,11]{x=y}{%
      \FormulaRefAuto{CaseFunInjectiveDisjointImagesCaseBB}{1,2,3,4,5,6,11}}
    \proofstepwidestar[1,2,3,4,5,6,7]{x=y}{\rOE{8,9,10,11,12}}
  \closeproofpartwideindR

  \proofpartwide{Schluss}
    \proofstepwidestar[1]{A\cap B=\varnothing}{\rA}
    \proofstepwidestar[2]{f\colon A\inj C}{\rA}
    \proofstepwidestar[3]{g\colon B\inj C}{\rA}
    \proofstepwidestar[4]{f[A]\cap g[B]=\varnothing}{\rA}
    \proofstepwidestar[2]{f\colon A\to C}{%
      \FormulaRefAuto{Injektive Funktion}{2}}
    \proofstepwidestar[3]{g\colon B\to C}{%
      \FormulaRefAuto{Injektive Funktion}{3}}
    \proofstepwidestar[2,3]{\CaseFun{A}{B}{f}{g}\colon A\cup B\to C}{%
      \FormulaRefAuto{f\colon A\to C\dsep g\colon B\to C
      \vdash \CaseFun{A}{B}{f}{g}\colon A\cup B\to C}{5,6}}
    \proofstepwidestar[8]{x\in A\cup B}{\rA}
    \proofstepwidestar[9]{y\in A\cup B}{\rA}
    \proofstepwidestar[10]{%
      \CaseFun{A}{B}{f}{g}(x)=\CaseFun{A}{B}{f}{g}(y)}{\rA}
    \proofstepwidestar[8]{x\in A\lor x\in B}{%
      \FormulaRefAuto{z \in A \cup B \eqvdash z \in A \lor z \in B}{8}}
    \proofstepwidestar[12]{x\in A}{\rA}
    \proofstepwidestar[1,2,3,4,9,10,12]{x=y}{%
      \FormulaRefAuto{CaseFunInjectiveDisjointImagesCaseA}{1,2,3,4,10,12,9}}
    \proofstepwidestar[14]{x\in B}{\rA}
    \proofstepwidestar[1,2,3,4,9,10,14]{x=y}{%
      \FormulaRefAuto{CaseFunInjectiveDisjointImagesCaseB}{1,2,3,4,10,14,9}}
    \proofstepwidestar[1,2,3,4,8,9,10]{x=y}{%
      \rOE{11,12,13,14,15}}
    \proofstepwidestar[1,2,3,4]{%
      \begin{aligned}[t]
      &\forall x,y\in A\cup B\,\\
      &\bigl(\CaseFun{A}{B}{f}{g}(x)=
      \CaseFun{A}{B}{f}{g}(y)\\
      &\rightarrow x=y\bigr)
      \end{aligned}}{\rUI{\rUI{\rRI{8,\rRI{9,\rRI{10,16}}}}}}
    \proofstepwidestar[1,2,3,4]{%
      \CaseFun{A}{B}{f}{g}\colon A\cup B\inj C}{%
      \FormulaRefAuto{Injektive Funktion}{7,17}}
  \closeproofpartwide
\end{tabproofsplitwide}

\FormulaThmDeltaK[Injektives Zusammenführen in disjunkte Zielmengen]{%
  A\cap B=\varnothing\dsep C\cap D=\varnothing\dsep
  f\colon A\inj C\dsep g\colon B\inj D
  \vdash \exists h\colon A\cup B\inj C\cup D
}{CaseFunDisjointTargetInjection}{%
  \DeltaRow{Mengen}{A\dsep B\dsep C\dsep D}
  \DeltaRow{Funktionen}{f\dsep g\dsep h}
}
\begin{tabproof}
  \proofstep{1}{A\cap B=\varnothing}{\rA}
  \proofstep{2}{C\cap D=\varnothing}{\rA}
  \proofstep{3}{f\colon A\inj C}{\rA}
  \proofstep{4}{g\colon B\inj D}{\rA}
  \proofstep{}{C\subseteq C\cup D}{\FormulaRefAuto{A\subseteq A\cup B}}
  \proofstep{}{D\subseteq C\cup D}{\FormulaRefAuto{B\subseteq A\cup B}}
  \proofstep{3}{f\colon A\to C}{\FormulaRefAuto{Injektive Funktion}{3}}
  \proofstep{4}{g\colon B\to D}{\FormulaRefAuto{Injektive Funktion}{4}}
  \proofstep{2,3,4}{f[A]\cap g[B]=\varnothing}{%
    \FormulaRefAuto{DisjointTargetImagesDisjoint}{7,8,2}}
  \proofstep{3}{f\colon A\inj C\cup D}{%
    \FormulaRefAuto{F\colon A\inj B\dsep B\subseteq C\vdash F\colon A\inj C}}
  \proofstep{4}{g\colon B\inj C\cup D}{%
    \FormulaRefAuto{F\colon A\inj B\dsep B\subseteq C\vdash F\colon A\inj C}}
  \proofstep{1,2,3,4}{\CaseFun{A}{B}{f}{g}\colon A\cup B\inj C\cup D}{%
    \FormulaRefAuto{CaseFunInjectiveDisjointImages}{1,10,11,9}}
  \proofstep{1,2,3,4}{\exists h\colon A\cup B\inj C\cup D}{\rEI{12}}
\end{tabproof}


\section{Inklusionsabbildung}

\subsection{Definition der Inklusionsabbildung}

\FormulaDefDeltaK[Inklusionsabbildung von \(A\) nach \(B\)]{%
  A\subseteq B\vdash
  \begin{aligned}[t]
    &\iota_{A,B}\colon A\to B,\\[-2pt]
    &\forall x\in A\,\bigl(\iota_{A,B}(x)\coloneqq x\bigr)
  \end{aligned}%
}{InclusionMapDef}{%
  \DeltaRow{Mengen}{A\dsep B}
  \DeltaRow{Funktionen}{\iota_{A,B}}
}
\begin{tabproof}
  \proofstep{1}{A\subseteq B}{\rA}
  \proofstep{2}{x\in A}{\rA}
  \proofstep{1,2}{x\in B}{%
    \FormulaRefAuto{A\subseteq B, x\in A \vdash x\in B}{1,2}}
\end{tabproof}

\FormulaThmDelta[Inklusionsabbildung von \(A\) nach \(B\)]
{A\subseteq B\vdash \iota_{A,B}\colon A\to B}
{
  \DeltaRow{Mengen}{A\dsep B}
}
\begin{tabproof}
  \proofstep{1}{A\subseteq B}{\rA}
  \proofstep{1}{\iota_{A,B}\colon A\to B}{%
    \FormulaRefAuto{InclusionMapDef}[def]{1}}
\end{tabproof}

\FormulaThmDelta
{A\subseteq B\vdash \TotRel{\iota_{A,B},A,B}}
{
  \DeltaRow{Mengen}{A\dsep B}
}
\begin{tabproof}
    \proofstep{1}{A\subseteq B}{\rA}
    \proofstep{1}{\iota_{A,B}\colon A\to B}{\FormulaRefAuto{A\subseteq B\vdash \iota_{A,B}\colon A\to B}{1}}
    \proofstep{1}{\TotRel{\iota_{A,B},A,B}}{\FormulaRefAuto{F\colon A\to B\vdash\TotRel{F,A,B}}{2}}
\end{tabproof}

\FormulaThmDelta[Grundgleichung der Inklusionsabbildung]
{A\subseteq B \dsep x\in A \vdash \iota_{A,B}(x)=x}
{
  \DeltaRow{Mengen}{A\dsep B\dsep x}
}
\begin{tabproof}
  \proofstep{1}{A\subseteq B}{\rA}
  \proofstep{2}{x\in A}{\rA}
  \proofstep{1}{\forall u\in A\,\bigl(\iota_{A,B}(u)=u\bigr)}{%
    \FormulaRefAuto{InclusionMapDef}[def]{1}}
  \proofstep{1,2}{\iota_{A,B}(x)=x}{%
    \rRE{\rUE{3},2}}
\end{tabproof}

\FormulaThmDelta
{A\subseteq B \dsep x\in A \vdash x=\iota_{A,B}(x)}%
{
  \DeltaRow{Mengen}{A\dsep B\dsep x}
}
\begin{tabproof}
  \proofstep{1}{A\subseteq B}{\rA}
  \proofstep{2}{x\in A}{\rA}

  \proofstep{1,2}{\iota_{A,B}(x)=x}{%
    \FormulaRefAuto{A\subseteq B \dsep x\in A \vdash \iota_{A,B}(x)=x}{1,2}}

  \proofstep{1,2}{x=\iota_{A,B}(x)}{%
    \FormulaRefAuto{a=b\vdash b=a}{3}}
\end{tabproof}


\subsection{Die Injektivität der Inklusionsabbildung}

\paragraph{Injektivität}
\FormulaThmDelta[Injektivität von \(\iota_{A,B}\)]{%
x\in A \dsep y\in A \dsep \iota_{A,B}(x)=\iota_{A,B}(y) \vdash x=y
}{
  \DeltaRow{Mengen}{A \dsep B \dsep x \dsep y}
}
\begin{tabproof}
  \proofstep{1}{\iota_{A,B}(x)=\iota_{A,B}(y)}{\rA}
  \proofstep{2}{x\in A}{\rA}
  \proofstep{3}{y\in A}{\rA}
  \proofstep{2}{\iota_{A,B}(x)=x}{%
    \FormulaRefAuto{A\subseteq B \dsep x\in A \vdash \iota_{A,B}(x)=x}{2}}
  \proofstep{3}{\iota_{A,B}(y)=y}{%
    \FormulaRefAuto{A\subseteq B \dsep x\in A \vdash \iota_{A,B}(x)=x}{3}}
  \proofstep{1,2,3}{x=y}{%
    \FormulaRefAuto{a=b\dsep a=c\dsep b=d\vdash c=d}{1,4,5}}
\end{tabproof}

\FormulaThmDelta[\(\iota_{A,B}\) als injektive Funktion]{%
A\subseteq B \vdash \iota_{A,B}\colon A\inj B
}{
  \DeltaRow{Mengen}{A\dsep B}
}
\begin{tabproof}
    \proofstep{1}{A\subseteq B}{\rA}
    \proofstep{1}{\iota_{A,B}\colon A\to B}{%
      \FormulaRefAuto{A\subseteq B\vdash \iota_{A,B}\colon A\to B}{1}}
    \proofstep{}{\forall x,y\in A\,(\iota_{A,B}(x)=\iota_{A,B}(y)\rightarrow x=y)}{%
      \FormulaRefAuto{x\in A \dsep y\in A \dsep \iota_{A,B}(x)=\iota_{A,B}(y) \vdash x=y}}
    \proofstep{1}{\iota_{A,B}\colon A\inj B}{%
      \FormulaRefAuto{Injektive Funktion}{2,3}}
\end{tabproof}



\section{Kompositionen}

\FormulaThmDelta{%
  B\subseteq C \dsep F\colon A\to B \vdash \iota_{B,C}\circ F\colon A\to C%
}{%
  \DeltaDecl{Mengen}{A\dsep B\dsep C}%
  \DeltaDecl{Funktionen}{F}%
}
\begin{tabproof}
  \proofstep{1}{B\subseteq C}{\rA}
  \proofstep{2}{F\colon A\to B}{\rA}
  \proofstep{1}{\iota_{B,C}\colon B\to C}{%
    \FormulaRefAuto{A\subseteq B\vdash \iota_{A,B}\colon A\to B}{1}}
  \proofstep{1,2}{\iota_{B,C}\circ F\colon A\to C}{%
    \FormulaRefAuto{G\circ F\colon A\to C}{2,3}}
\end{tabproof}

\subsection{Inklusionsabbildung und Komposition}

\FormulaThmDelta[Werte der Inklusionskomposition liegen im kleineren Ziel]{%
  A\subseteq B\dsep F\colon C\to A
  \vdash
  \forall x\in C\,(\iota_{A,B}\circ F)(x)\in A
}{%
  \DeltaRow{Mengen}{A\dsep B\dsep C\dsep x}
  \DeltaRow{Funktionen}{F}
}
\begin{tabproof}
  \proofstep{1}{A\subseteq B}{\rA}
  \proofstep{2}{F\colon C\to A}{\rA}
  \proofstep{3}{x\in C}{\rA}

  \proofstep{2,3}{F(x)\in A}{%
    \FormulaRefAuto{F\colon A\to B\dsep x\in A\vdash F(x)\in B}{2,3}}
  \proofstep{3}{(\iota_{A,B}\circ F)(x)=\iota_{A,B}(F(x))}{%
    \FormulaRefAuto{x\in A \vdash (F\circ G)(x)=F(G(x))}{3}}
  \proofstep{1,2,3}{\iota_{A,B}(F(x))=F(x)}{%
    \FormulaRefAuto{A\subseteq B \dsep x\in A \vdash \iota_{A,B}(x)=x}{1,4}}
  \proofstep{1,2,3}{(\iota_{A,B}\circ F)(x)=F(x)}{%
    \FormulaRefAuto{a=b,\, b=c \vdash a=c}{5,6}}
  \proofstep{1,2,3}{F(x)=(\iota_{A,B}\circ F)(x)}{%
    \FormulaRefAuto{a=b\vdash b=a}{7}}
  \proofstep{1,2,3}{(\iota_{A,B}\circ F)(x)\in A}{\rIE{8,4}}
  \proofstep{1,2}{\forall x\in C\,(\iota_{A,B}\circ F)(x)\in A}{%
    \rUI{\rRI{3,9}}}
\end{tabproof}

\FormulaThmDeltaK[Auswertung einer Inklusionskomposition]{%
  B\subseteq C\dsep F\colon A\to B\dsep x\in A
  \vdash (\iota_{B,C}\circ F)(x)=F(x)%
}{InclusionCompositionValue}{%
  \DeltaRow{Mengen}{A\dsep B\dsep C\dsep x}%
  \DeltaRow{Funktionen}{F}%
}
\begin{tabproof}
  \proofstep{1}{B\subseteq C}{\rA}
  \proofstep{2}{F\colon A\to B}{\rA}
  \proofstep{3}{x\in A}{\rA}
  \proofstep{2,3}{F(x)\in B}{%
    \FormulaRefAuto{F\colon A\to B\dsep x\in A\vdash F(x)\in B}{2,3}}
  \proofstep{3}{(\iota_{B,C}\circ F)(x)=\iota_{B,C}(F(x))}{%
    \FormulaRefAuto{x\in A \vdash (F\circ G)(x)=F(G(x))}{3}}
  \proofstep{1,2,3}{\iota_{B,C}(F(x))=F(x)}{%
    \FormulaRefAuto{A\subseteq B \dsep x\in A \vdash \iota_{A,B}(x)=x}{1,4}}
  \proofstep{1,2,3}{(\iota_{B,C}\circ F)(x)=F(x)}{%
    \FormulaRefAuto{a=b,\,b=c\vdash a=c}{5,6}}
\end{tabproof}


\subsection{Erhalt der Injektivität}

% — Erhalt von Injektivität / Surjektivität / Bijektivität —
\FormulaThmDelta[Erhalt der Injektivität]{%
x\in A \dsep y\in A \dsep (G\circ F)(x)=(G\circ F)(y) \vdash x=y
}{
  \DeltaRow{Mengen}{x \dsep y \dsep A \dsep B \dsep C}
  \DeltaRow{Injektive Funktionen}{F\colon A\inj B \dsep G\colon B\inj C}
}
\begin{tabproof}
  \proofstep{1}{(G\circ F)(x)=(G\circ F)(y)}{\rA}
  \proofstep{2}{x\in A}{\rA}
  \proofstep{3}{y\in A}{\rA}
  \proofstep{2}{F(x)\in B}{\FormulaRefAuto{F\colon A\to B\dsep x\in A\vdash F(x)\in B}{2}}
  \proofstep{3}{F(y)\in B}{\FormulaRefAuto{F\colon A\to B\dsep x\in A\vdash F(x)\in B}{3}}
  \proofstep{1}{G(F(x))=G(F(y))}{\rIE{\FormulaRefAuto{x\in A \vdash (G\circ F)(x)=G(F(x))},1}}
  \proofstep{1,2,3}{F(x)=F(y)}{\FormulaRefAuto{Injektivität}{4,5,6}}
  \proofstep{1,2,3}{x=y}{\FormulaRefAuto{Injektivität}{2,3,7}}
\end{tabproof}


\FormulaThmDeltaR[Die Komposition als injektive Funktion]{%
G\circ F\colon A\inj C
}{F\colon A\inj B \dsep G\colon B\inj C\vdash G\circ F\colon A\inj C}{
  \DeltaRow{Mengen}{A\dsep B\dsep C}
  \DeltaPrem{Injektive Funktionen}{F\colon A\inj B \dsep G\colon B\inj C}
}
\begin{tabproof}
    \proofstep{}{G\circ F\colon A\to C}{\FormulaRefAuto{G\circ F\colon A\to C}}
    \proofstep{}{\forall x,y\in A\,((G\circ F)(x)=(G\circ F)(y) \rightarrow x=y)}{\FormulaRefAuto{x\in A \dsep y\in A \dsep (G\circ F)(x)=(G\circ F)(y) \vdash x=y}}
    \proofstep{}{G\circ F\colon A\inj C}{\FormulaRefAuto{Injektive Funktion}{1,2}}
\end{tabproof}

\FormulaThmDelta{%
  B\subseteq C \dsep F\colon A\inj B \vdash \iota_{B,C}\circ F\colon A\inj C%
}{%
  \DeltaDecl{Mengen}{A\dsep B\dsep C}%
  \DeltaDecl{Funktionen}{F}%
}
\begin{tabproof}
  \proofstep{1}{B\subseteq C}{\rA}
  \proofstep{2}{F\colon A\inj B}{\rA}
  \proofstep{1}{\iota_{B,C}\colon B\inj C}{%
    \FormulaRefAuto{A\subseteq B\vdash \iota_{A,B}\colon A\inj B}{1}}
  \proofstep{1,2}{\iota_{B,C}\circ F\colon A\inj C}{%
    \FormulaRefAuto{F\colon A\inj B \dsep G\colon B\inj C\vdash G\circ F\colon A\inj C}{2,3}}
\end{tabproof}


\subsection{Komposition mit Inklusionen und Einschränkung}

\FormulaThmDelta{%
A\subseteq B\vdash F\restriction_A = F\circ \iota_{A,B}%
}{
  \DeltaRow{Mengen}{A \dsep B \dsep C \dsep x}
  \DeltaPrem{Funktionen}{F\colon B\to C}
  \DeltaPrem{Inklusionsabbildung}{\iota_{A,B}\colon A\to B}[\FormulaRefAuto{A\subseteq B\vdash \iota_{A,B}\colon A\to B}]
  \DeltaPrem{Komposition}{F\circ \iota_{A,B}\colon A\to C}[\FormulaRefAuto{G\circ F\colon A\to C}]
  \DeltaPrem{Einschränkung}{F\restriction_A\colon A\to C}[\FormulaRefAuto{C\subseteq A \vdash F\restriction_C\colon C\to B}]
}
\begin{tabproofwide}
  \proofstepwidestar[1]{A\subseteq B}{\rA}

  \proofstepwidestar[2]{x\in A}{\rA}

  \proofstepwidestar[1,2]{x=\iota_{A,B}(x)}{%
    \FormulaRefAuto{A\subseteq B \dsep x\in A \vdash x=\iota_{A,B}(x)}{1,2}}

  \proofstepwide[1,2]{F\restriction_A(x)}{=}{F(x)}{
    \FormulaRefAuto{C\subseteq A \dsep x\in C \vdash F\restriction_C(x)=F(x)}{1,2}}

  \proofstepwide[1,2]{}{=}{F(\iota_{A,B}(x))}{
    \rIE{3,4}}

  \proofstepwide[1,2]{}{=}{(F\circ\iota_{A,B})(x)}{
    \FormulaRefAuto{x\in A \vdash F(G(x))=(F\circ G)(x)}{2}}

  \proofstepwide[1,2]{F\restriction_A(x)}{=}{(F\circ\iota_{A,B})(x)}{
    \rChain{4,6}}

  \proofstepwidestar[1]{\forall x\in A\,\bigl(F\restriction_A(x)=(F\circ \iota_{A,B})(x)\bigr)}{%
    \rUI{\rRI{2,7}}}

  \proofstepwidestar[1]{F\restriction_A\colon A\to C}{%
    \FormulaRefAuto{C\subseteq A \vdash F\restriction_C\colon C\to B}{1}}

  \proofstepwidestar[1]{F\circ\iota_{A,B}\colon A\to C}{%
    \FormulaRefAuto{G\circ F\colon A\to C}}

  \proofstepwidestar[1]{F\restriction_A = F\circ \iota_{A,B}}{%
    \FormulaRefAuto{FunctionExtensionality}[thm]{9,10,8}}
\end{tabproofwide}


\section{Leere Bereiche}

\FormulaThmDelta[Leere Relation ist injektiv]{%
  \varnothing\colon \varnothing\inj M%
}{%
  \DeltaDecl{Mengen}{M\dsep x\dsep y}%
}
\begin{tabproof}
  \proofstep{}{ \varnothing\colon \varnothing\to M }{%
    \FormulaRefAuto{\varnothing\colon \varnothing\to M}}

  \proofstep{2}{x\in \varnothing \land y\in \varnothing \land \varnothing(x)=\varnothing(y)}{\rA}
  \proofstep{2}{x\in \varnothing}{\rAEa{2}}
  \proofstep{}{x\notin \varnothing}{\FormulaRefAuto{x \not\in \varnothing}}
  \proofstep{2}{\bot}{\rBI{3,4}}
  \proofstep{2}{x=y}{\FormulaRefAuto{P\dsep \neg P \vdash Q}{3,4}}
  \proofstep{}{(x\in \varnothing \land y\in \varnothing \land \varnothing(x)=\varnothing(y))\rightarrow x=y}{\rRI{2,6}}

  \proofstep{}{ \varnothing\colon \varnothing\inj M }{%
    \FormulaRefAuto{Injektive Funktion}{1,7}}
\end{tabproof}


\section{Adjunktionsabbildung einer Mengenfamilie}

Als abschließende Konstruktion dieses Kapitels betrachten wir eine spezielle
injektive Abbildung auf Mengenfamilien. Sie steht bereits hier statt erst im
Frankl-Band, weil sie nur den allgemeinen Funktionsbegriff aus Band~05 und den
Injektivitätsbegriff dieses Bandes voraussetzt. Die Stellung nach den leeren
Bereichen ist zugleich sachlich passend: Im Gegenbeispiel zur Surjektivität
tritt ein leerer Definitionsbereich auf.

Für eine Familie \(M\) und ein Element \(a\) zerlegen die folgenden beiden
Teilmengenfamilien \(M\) danach, ob ihre Elemente \(a\) enthalten oder
vermeiden. Sie bilden den natürlichen Ziel- und Definitionsbereich der
punktweisen Adjunktion \(X\mapsto X\cup\{a\}\).

\FormulaDefDeltaK[Enthaltende Teilfamilie zu einem Element]{%
  \FamContaining{M}{a} \coloneqq \{X\in M\mid a\in X\}
}{Enthaltende Teilfamilie}{%
  \DeltaRow{Mengen}{M\dsep X\dsep a}
  \DeltaRow{Neue Symbole}{\FamContaining{M}{a}}
}

\FormulaDefDeltaK[Vermeidende Teilfamilie zu einem Element]{%
  \FamAvoiding{M}{a} \coloneqq \{X\in M\mid a\notin X\}
}{Vermeidende Teilfamilie}{%
  \DeltaRow{Mengen}{M\dsep X\dsep a}
  \DeltaRow{Neue Symbole}{\FamAvoiding{M}{a}}
}

Unter Vereinigungsschließung und \(\{a\}\in M\) wird die Adjunktion zu einer
Funktion zwischen diesen Teilfamilien.

\FormulaDefDeltaK[Adjunktionsabbildung zur Einermenge]{%
  \UnionClosedFam(M)\dsep \{a\}\in M\vdash
  \begin{aligned}[t]
    &\FranklAdj{M}{a}\colon
      \FamAvoiding{M}{a}\to \FamContaining{M}{a},\\[-2pt]
    &\forall X\in \FamAvoiding{M}{a}\,
      \bigl(\FranklAdj{M}{a}(X)\coloneqq X\cup\{a\}\bigr)
  \end{aligned}
}{Adjunktionsabbildung einer Mengenfamilie}{%
  \DeltaRow{Mengen}{M\dsep X\dsep Y\dsep a}
  \DeltaRow{Funktionen}{\FranklAdj{M}{a}}
  \DeltaRow{Teilmengenfamilien}{\FamContaining{M}{a}\dsep
    \FamAvoiding{M}{a}}
  \DeltaRow{Neue Symbole}{\FranklAdj{M}{a}}
}
\begin{tabproof}
  \proofstep{1}{\UnionClosedFam(M)}{\rA}
  \proofstep{2}{\{a\}\in M}{\rA}
  \proofstep{3}{X\in \FamAvoiding{M}{a}}{\rA}
  \proofstep{3}{X\in M\land a\notin X}{%
    \FormulaRefAuto{\FamAvoiding{M}{a}\coloneqq
      \{X\in M\mid a\notin X\}}[def]{3}}
  \proofstep{3}{X\in M}{\rAEa{4}}
  \proofstep{1,2,3}{X\cup\{a\}\in M}{%
    \FormulaRefAuto{Vereinigungsabgeschlossene Familie}[def]{1,5,2}}
  \proofstep{}{a\in X\cup\{a\}}{%
    \FormulaRefAuto{a\in A\cup\{a\}}[thm]}
  \proofstep{1,2,3}{X\cup\{a\}\in M\land a\in X\cup\{a\}}{%
    \rAI{6,7}}
  \proofstep{1,2,3}{X\cup\{a\}\in \FamContaining{M}{a}}{%
    \FormulaRefAuto{\FamContaining{M}{a}\coloneqq
      \{X\in M\mid a\in X\}}[def]{8}}
\end{tabproof}

\FormulaThmDelta[Typisierung der Adjunktionsabbildung]{%
  \UnionClosedFam(M)\dsep \{a\}\in M
  \vdash \FranklAdj{M}{a}\colon
  \FamAvoiding{M}{a}\to \FamContaining{M}{a}
}{%
  \DeltaRow{Mengen}{M\dsep X\dsep a}
  \DeltaRow{Funktionen}{\FranklAdj{M}{a}}
}
\begin{tabproof}
  \proofstep{1}{\UnionClosedFam(M)}{\rA}
  \proofstep{2}{\{a\}\in M}{\rA}
  \proofstep{1,2}{\FranklAdj{M}{a}\colon
    \FamAvoiding{M}{a}\to \FamContaining{M}{a}}{%
    \FormulaRefAuto{Adjunktionsabbildung einer Mengenfamilie}[def]{1,2}}
\end{tabproof}

\FormulaThmDelta[Explizite Werte der Adjunktionsabbildung]{%
  \UnionClosedFam(M)\dsep \{a\}\in M\dsep X\in \FamAvoiding{M}{a}
  \vdash \FranklAdj{M}{a}(X)=X\cup\{a\}
}{%
  \DeltaRow{Mengen}{M\dsep X\dsep a}
  \DeltaRow{Funktionen}{\FranklAdj{M}{a}}
}
\begin{tabproof}
  \proofstep{1}{\UnionClosedFam(M)}{\rA}
  \proofstep{2}{\{a\}\in M}{\rA}
  \proofstep{3}{X\in \FamAvoiding{M}{a}}{\rA}
  \proofstep{1,2}{%
    \forall Y\in \FamAvoiding{M}{a}\,
      \bigl(\FranklAdj{M}{a}(Y)=Y\cup\{a\}\bigr)}{%
    \FormulaRefAuto{Adjunktionsabbildung einer Mengenfamilie}[def]{1,2}}
  \proofstep{1,2,3}{\FranklAdj{M}{a}(X)=X\cup\{a\}}{%
    \rRE{\rUE{4},3}}
\end{tabproof}

\FormulaThmDelta[Adjunktionsabbildung ist injektiv]{%
  \UnionClosedFam(M)\dsep \{a\}\in M
  \vdash \FranklAdj{M}{a}\colon
  \FamAvoiding{M}{a}\inj \FamContaining{M}{a}
}{%
  \DeltaRow{Mengen}{M\dsep X\dsep Y\dsep a}
  \DeltaRow{Funktionen}{\FranklAdj{M}{a}}
}
\begin{tabproof}
  \proofstep{1}{\UnionClosedFam(M)}{\rA}
  \proofstep{2}{\{a\}\in M}{\rA}
  \proofstep{1,2}{\FranklAdj{M}{a}\colon
    \FamAvoiding{M}{a}\to \FamContaining{M}{a}}{%
    \FormulaRefAuto{\UnionClosedFam(M)\dsep \{a\}\in M
    \vdash \FranklAdj{M}{a}\colon
    \FamAvoiding{M}{a}\to \FamContaining{M}{a}}{1,2}}

  \proofstep{4}{X\in \FamAvoiding{M}{a}}{\rA}
  \proofstep{5}{Y\in \FamAvoiding{M}{a}}{\rA}
  \proofstep{6}{\FranklAdj{M}{a}(X)=\FranklAdj{M}{a}(Y)}{\rA}
  \proofstep{1,2,4}{\FranklAdj{M}{a}(X)=X\cup\{a\}}{%
    \FormulaRefAuto{\UnionClosedFam(M)\dsep \{a\}\in M\dsep
    X\in \FamAvoiding{M}{a}\vdash
    \FranklAdj{M}{a}(X)=X\cup\{a\}}{1,2,4}}
  \proofstep{1,2,5}{\FranklAdj{M}{a}(Y)=Y\cup\{a\}}{%
    \FormulaRefAuto{\UnionClosedFam(M)\dsep \{a\}\in M\dsep
    X\in \FamAvoiding{M}{a}\vdash
    \FranklAdj{M}{a}(X)=X\cup\{a\}}{1,2,5}}
  \proofstep{1,2,4}{X\cup\{a\}=\FranklAdj{M}{a}(X)}{%
    \FormulaRefAuto{a=b\vdash b=a}{7}}
  \proofstep{1,2,4,6}{X\cup\{a\}=\FranklAdj{M}{a}(Y)}{%
    \FormulaRefAuto{a = b,\, b = c \vdash a = c}{9,6}}
  \proofstep{1,2,4,5,6}{X\cup\{a\}=Y\cup\{a\}}{%
    \FormulaRefAuto{a = b,\, b = c \vdash a = c}{10,8}}
  \proofstep{4}{a\notin X}{%
    \rAEb{\FormulaRefAuto{\FamAvoiding{M}{a}\coloneqq
    \{X\in M\mid a\notin X\}}{4}}}
  \proofstep{5}{a\notin Y}{%
    \rAEb{\FormulaRefAuto{\FamAvoiding{M}{a}\coloneqq
    \{X\in M\mid a\notin X\}}{5}}}
  \proofstep{1,2,4,5,6}{X=Y}{%
    \FormulaRefAuto{a\notin X\dsep a\notin Y\dsep
    X\cup\{a\}=Y\cup\{a\}\vdash X=Y}{12,13,11}}

  \proofstep{1,2}{\FranklAdj{M}{a}\colon
    \FamAvoiding{M}{a}\inj \FamContaining{M}{a}}{%
    \FormulaRefAuto{Injektive Funktion}{%
      3,\rUI{\rRI{4,\rRI{5,\rRI{6,14}}}}}}
\end{tabproof}

Die Adjunktionsabbildung ist im Allgemeinen nicht surjektiv. Für
\(M=\{\{a\},\{a,b\}\}\) mit \(a\neq b\) ist \(M\)
vereinigungsabgeschlossen und enthält \(\{a\}\), aber
\(\FamAvoiding{M}{a}=\varnothing\) und
\(\FamContaining{M}{a}=M\neq\varnothing\). Daher ist die Injektivität ihre
kennzeichnende allgemeine Abbildungseigenschaft; eine Bijektion liegt ohne
weitere Voraussetzungen ebenfalls nicht vor.

\chapter{Weitere Anwendungen}

\section{Weitere injektive Größenvergleiche}

\FormulaDefDeltaK[Injektive Form von höchstens der Hälfte]{%
  B\subseteq A\vdash
  \AtMostHalf{B}{A}\coloneqq
  \exists F\colon B\inj A\setminus B
}{AtMostHalf}{%
  \DeltaRow{Mengen}{A\dsep B}
  \DeltaRow{Funktionen}{F}
  \DeltaRow{\textbf{Neues Symbol}}{}
  \DeltaPrem{Höchstens die Hälfte}{\AtMostHalf{B}{A}}
}

\FormulaThmDeltaK[Eine Einermenge injiziert in jede nichtleere Menge]{%
  b\in B\vdash \exists F\colon\{a\}\inj B
}{SingletonInjectionIntoNonempty}{%
  \DeltaRow{Mengen}{A\dsep B\dsep a\dsep b\dsep x\dsep y}
  \DeltaRow{Funktionen}{F\dsep\ConstFun{\{a\}}{B}{b}}
}
\begin{tabproof}
  \proofstep{1}{b\in B}{\rA}
  \proofstep{1}{\ConstFun{\{a\}}{B}{b}\colon\{a\}\to B}{%
    \FormulaRefAuto{b\in B\vdash \ConstFun{A}{B}{b}\colon A\to B}{1}}
  \proofstep{3}{x\in\{a\}}{\rA}
  \proofstep{4}{y\in\{a\}}{\rA}
  \proofstep{5}{\ConstFun{\{a\}}{B}{b}(x)=
    \ConstFun{\{a\}}{B}{b}(y)}{\rA}
  \proofstep{3}{x=a}{\FormulaRefAuto{x\in\{a\}\eqvdash x=a}{3}}
  \proofstep{4}{y=a}{\FormulaRefAuto{x\in\{a\}\eqvdash x=a}{4}}
  \proofstep{3,4}{x=y}{%
    \FormulaRefAuto{a=b,\,c=b\vdash a=c}{6,7}}
  \proofstep{1}{\ConstFun{\{a\}}{B}{b}\colon\{a\}\inj B}{%
    \FormulaRefAuto{Injektive Funktion}[def]{2,3,4,5,8}}
  \proofstep{1}{\exists F\colon\{a\}\inj B}{\rEI{9}}
\end{tabproof}

\end{document}
